\documentclass[12pt,a4paper]{article}

% Encoding and fonts — LuaLaTeX/XeLaTeX
\usepackage{fontspec}
\setmainfont{FreeSerif}
\setmonofont{FreeMono}

% Page layout
\usepackage[top=1in, bottom=1in, left=1in, right=1in]{geometry}

% Typography
\usepackage{microtype}
\usepackage{parskip}

% Language
\usepackage[english]{babel}

% Headers and footers
\usepackage{fancyhdr}
\pagestyle{fancy}
\fancyhf{}
\fancyhead[L]{\leftmark}
\fancyhead[R]{\thepage}
\fancyfoot[C]{\thepage}
\renewcommand{\headrulewidth}{0.4pt}
\renewcommand{\footrulewidth}{0pt}

% Section styling
\usepackage{titlesec}
\titleformat{\section}{\Large\bfseries}{\thesection}{1em}{}
\titleformat{\subsection}{\large\bfseries}{\thesubsection}{1em}{}
\titleformat{\subsubsection}{\normalsize\bfseries}{\thesubsubsection}{1em}{}
\titlespacing*{\section}{0pt}{2.5ex plus 1ex minus .2ex}{1.5ex plus .2ex}
\titlespacing*{\subsection}{0pt}{2ex plus 1ex minus .2ex}{1ex plus .2ex}
\titlespacing*{\subsubsection}{0pt}{1.5ex plus 1ex minus .2ex}{0.8ex plus .2ex}

% List formatting
\usepackage{enumitem}
\setlist[itemize]{topsep=0.5ex, itemsep=0.25ex, parsep=0pt}
\setlist[enumerate]{topsep=0.5ex, itemsep=0.25ex, parsep=0pt}

% Essential packages
\usepackage{graphicx}
\usepackage[table]{xcolor}
\usepackage{booktabs}
\usepackage{array}
\usepackage{tabularx}
\usepackage{longtable}
\usepackage{amsmath}
\usepackage{amssymb}
\usepackage[normalem]{ulem}
\renewcommand{\arraystretch}{1.3}

% Cross-references
\usepackage{hyperref}
\usepackage{cleveref}

% Custom commands
\newcommand{\pilcrow}{\S}

% Hyperref setup
\hypersetup{
    colorlinks=true,
    linkcolor=blue,
    filecolor=magenta,
    urlcolor=cyan,
}

% Code listings
\usepackage{listings}
\definecolor{codebg}{RGB}{248,248,248}
\definecolor{codeframe}{RGB}{200,200,200}
\definecolor{codekeyword}{RGB}{0,0,180}
\definecolor{codecomment}{RGB}{0,128,0}
\definecolor{codestring}{RGB}{163,21,21}
\lstset{
    basicstyle=\ttfamily\small,
    breaklines=true,
    breakatwhitespace=false,
    frame=single,
    framerule=0.5pt,
    rulecolor=\color{codeframe},
    backgroundcolor=\color{codebg},
    keepspaces=true,
    numbers=left,
    numberstyle=\tiny\color{gray},
    numbersep=8pt,
    showstringspaces=false,
    tabsize=4,
    xleftmargin=1em,
    xrightmargin=1em,
    aboveskip=1em,
    belowskip=1em,
    keywordstyle=\color{codekeyword}\bfseries,
    commentstyle=\color{codecomment}\itshape,
    stringstyle=\color{codestring},
    literate={_}{{\textunderscore}}1
             {-}{{\textminus}}1
             {<}{{\textless}}1
             {>}{{\textgreater}}1
             {~}{{\textasciitilde}}1
             {^}{{\textasciicircum}}1
}

% YAML language definition for listings
\lstdefinelanguage{YAML}{
    keywords={true,false,null,y,n},
    keywordstyle=\color{blue}\bfseries,
    sensitive=false,
    comment=[l]{\#},
    morecomment=[s]{/*}{*/},
    commentstyle=\color{gray}\ttfamily,
    stringstyle=\color{red}\ttfamily,
    morestring=[b]',
    morestring=[b]',
}

% TOML language definition for listings
\lstdefinelanguage{TOML}{
    keywords={true,false},
    keywordstyle=\color{blue}\bfseries,
    sensitive=true,
    comment=[l]{\#},
    commentstyle=\color{gray}\ttfamily,
    stringstyle=\color{red}\ttfamily,
    morestring=[b]',
    morestring=[b]',
}

% Dockerfile language definition for listings
\lstdefinelanguage{Dockerfile}{
    keywords={FROM,RUN,CMD,LABEL,EXPOSE,ENV,ADD,COPY,ENTRYPOINT,VOLUME,USER,WORKDIR,ARG,ONBUILD,STOPSIGNAL,HEALTHCHECK,SHELL},
    keywordstyle=\color{blue}\bfseries,
    sensitive=false,
    comment=[l]{\#},
    commentstyle=\color{gray}\ttfamily,
    stringstyle=\color{red}\ttfamily,
    morestring=[b]',
    morestring=[b]',
}

% JSON language definition for listings
\lstdefinelanguage{JSON}{
    keywords={true,false,null},
    keywordstyle=\color{blue}\bfseries,
    sensitive=true,
    commentstyle=\color{gray}\ttfamily,
    stringstyle=\color{red}\ttfamily,
    morestring=[b]',
}

% Code listing float environment
\usepackage{float}
\newfloat{listing}{htbp}{lol}[section]
\floatname{listing}{Listing}

% Admonition boxes
\usepackage{tcolorbox}
\tcbuselibrary{skins,breakable}

% Admonition style definitions
\newtcolorbox{admonitionbox}[2][]{
    enhanced,
    breakable,
    colback=#2!5,
    colframe=#2!80!black,
    fonttitle=\bfseries,
    title=#1,
    left=2mm,
    right=2mm,
}

% Synthon tuple boxes
\newtcolorbox{synthonbox}{
    enhanced,
    colback=blue!5,
    colframe=blue!75!black,
    fontupper=\large,
    center upper,
    boxrule=1pt,
    arc=4pt,
}

\begin{document}

\begin{center}
{\Large\bfseries ZFC$_t$ Functor Discovery: Overview of Results}

\rule{0.5\textwidth}{0.4pt}
\end{center}

\subsection*{What Was Built}

The manipulator is a functor discovery machine. It applies algebraic operations ---
tensor product (per-primitive supremum), meet (per-primitive infimum), and single-primitive lift ---
to imscription tuples, and exposes the corresponding ZFC$_t$ clause transformations side by side.

The intent is to reverse-engineer the composition rules that govern how ZFC$_t$ expressions combine, and to identify where non-trivial cross-clause phenomena emerge that are not derivable from the per-primitive rules alone.

\begin{center}
\rule{0.5\textwidth}{0.4pt}
\end{center}

\subsection*{Tier Landscape of the Special Entries}

The nine canonical reference entries divide cleanly across three tier levels.

\subsubsection*{$O_0$ --- Semasiographic baseline ($\hat{\varphi} =$ non-critical)}

\begin{table}[htbp]
\centering
\small

\begin{tabularx}{\textwidth}{>{\raggedright\arraybackslash}X >{\centering\arraybackslash}X >{\centering\arraybackslash}X >{\raggedright\arraybackslash}X}
\toprule
\textbf{Entry} & \textbf{FROB} & \textbf{FIXPT} & \textbf{Notable} \\
\midrule

ZFC\_foundations & yes & no & Has Frobenius ($\Phi_{\}}$, ord 4) --- but no criticality \\
heat\_diffusion\_equation & no & no & SEQAX present; no criticality \\
wave\_equation\_temporal & yes & no & Has Frobenius; no criticality \\
\bottomrule
\end{tabularx}
\end{table}

The presence of Frobenius ($\Phi_{\}}$) does \textbf{not} lift an entry out of $O_0$. \\
Tier assignment is gated first on criticality ($\hat{\varphi}$ ordinal $\geq 1$). Without it, the Frobenius structure is algebraically present but ourobourically inert.

\subsubsection*{$O_2^\dagger$ --- Recursive, $D_\infty$, below Frobenius cliff}

\begin{table}[htbp]
\centering
\small

\begin{tabularx}{\textwidth}{>{\raggedright\arraybackslash}X >{\centering\arraybackslash}X >{\centering\arraybackslash}X >{\raggedright\arraybackslash}X}
\toprule
\textbf{Entry} & \textbf{FROB} & \textbf{FIXPT} & \textbf{Promo atoms} \\
\midrule

ZFCt & no & yes & HOLOBOUND LR\_DUAL PM\_Z2 SEQAX TEMPD2 ZWIND \\
Schrödinger equation & no & yes & LR\_DUAL SEQAX TEMPD2 ZWIND \\
Navier-Stokes equations & no & yes & LR\_DUAL PM\_Z2 SEQAX TEMPD2 ZWIND \\
\bottomrule
\end{tabularx}
\end{table}

All three are critical ($\hat{\varphi}_{\text{ÿ}}$, FIXPT), temporally extended ($\Omega_{\text{z}}$, $D_\infty$), but lack Frobenius symmetry ($\Phi <$ ord 4). They sit immediately below the cliff.

\subsubsection*{$O_\infty$ --- Frobenius round-trip}

\begin{table}[htbp]
\centering
\small

\begin{tabularx}{\textwidth}{>{\raggedright\arraybackslash}X >{\centering\arraybackslash}X >{\centering\arraybackslash}X >{\raggedright\arraybackslash}X}
\toprule
\textbf{Entry} & \textbf{FROB} & \textbf{FIXPT} & \textbf{Promo atoms} \\
\midrule

Temporal Mathematics & yes & yes & HOLOBOUND SEQAX TEMPD2 ZWIND \\
Einstein field equations & yes & yes & HOLOBOUND SEQAX TEMPD2 ZWIND \\
IUG (Mochizuki) & yes & yes & HOLOBOUND LR\_DUAL SEQAX ZWIND \\
\bottomrule
\end{tabularx}
\end{table}

All three independently satisfy the Frobenius round-trip condition (R1): \\
$\hat{\varphi} \geq$ ord 1 and $\Phi = \Phi_{\}}$. \\
FROB and FIXPT are both present in each.

\begin{center}
\rule{0.5\textwidth}{0.4pt}
\end{center}

\subsection*{The Central Result: $\mathrm{tensor}(\mathrm{ZFC},\,\mathrm{ZFCt}) = O_\infty$}

\begin{lstlisting}
ZFC     tier = O0      Phi = Phi_}  (ord 4, FROB)    phi^ = phi^_z  (ord 0, no FIXPT)
ZFCt    tier = O2+     Phi = Phi_F  (ord 2, no FROB) phi^ = phi^_y  (ord 1, FIXPT)
--------------------------------------------------------------------
tensor  tier = O_inf   Phi = Phi_}  (from ZFC)       phi^ = phi^_y  (from ZFCt)
\end{lstlisting}

ZFC holds the Frobenius component. ZFCt holds the criticality component. \\
Neither alone satisfies R1 (FROB $\land$ FIXPT). Their tensor product satisfies both simultaneously --- tier emergence to $O_\infty$.

\textbf{Per-primitive breakdown:} \\
The tensor sources eleven of twelve primitives from ZFCt (all six promotion channels plus ƒ, $\Gamma$, $\Sigma$, Ħ, $\Omega$), and exactly one from ZFC: $\Phi$. \\
ZFCt's $\Phi_{\text{F}}$ (ord 2) is lower than ZFC's $\Phi_{\}}$ (ord 4), so the supremum retains ZFC's Frobenius.

\begin{lstlisting}
d(ZFC, ZFCt)    = 7.148
d(ZFC, tensor)  = 6.863   --- tensor is almost as far from ZFC as ZFCt is
d(ZFCt, tensor) = 2.000   --- tensor is just one ordinal step from ZFCt (only Phi differs)
\end{lstlisting}

The meet falls to $O_0$. The infimum is destructive --- it strips the temporal promotions and recovers the floor.

\subsubsection*{The Asymmetry in $\Phi$}

The six ZFC$_t$ promotion channels all move upward in ordinal from ZFC base values:

\begin{table}[htbp]
\centering
\small

\begin{tabularx}{\textwidth}{>{\centering\arraybackslash}X >{\centering\arraybackslash}X >{\centering\arraybackslash}X >{\centering\arraybackslash}X}
\toprule
\textbf{Primitive} & \textbf{ZFC base} & \textbf{$\rightarrow$ ZFCt promoted} & \textbf{Gap} \\
\midrule

Þ & $Þ_{\text{6}}$ (ord 0) & $Þ_{\text{O}}$ (ord 4) & +4 \\
Ř & $Ř_{\text{¯}}$ (ord 0) & $Ř_{\text{=}}$ (ord 3) & +3 \\
ɢ & $ɢ_{\text{^}}$ (ord 0) & $ɢ_{\text{ˌ}}$ (ord 2) & +2 \\
Ħ & $Ħ_{\text{Ñ}}$ (ord 0) & $Ħ_{\text{A}}$ (ord 2) & +2 \\
$\Omega$ & $\Omega_{\text{Å}}$ (ord 0) & $\Omega_{\text{z}}$ (ord 2) & +2 \\
$\Phi$ & $\Phi_{\text{ɐ}}$ (ord 0) & $\Phi_{\text{F}}$ (ord 2) & +2 \\
\bottomrule
\end{tabularx}
\end{table}

ZFCt's promotion for $\Phi$ targets $\Phi_{\text{F}}$ ($\pm$-symmetry, ord 2). \\
However, the actual \texttt{ZFC\_foundations} entry carries $\Phi_{\text{˙}}$ (normalized to $\Phi_{\}}$, ord 4) --- which sits \textbf{above} the ZFCt-promoted value. \\
Thus \textbf{ZFCt is strictly lower than ZFC on the $\Phi$ axis.}

ZFCt trades Frobenius for temporal-sequential structure. It acquires six new capabilities (holographic topology, lateral relations, $\mathbb{Z}_2$ parity, sequentiality, chirality, integer winding) but deliberately operates below the Frobenius cliff. \\
ZFC carries the Frobenius symmetry but is non-critical and non-temporal.

They are \textbf{complementary}, not hierarchical. Their tensor product is their reunion.

\begin{center}
\rule{0.5\textwidth}{0.4pt}
\end{center}

\subsection*{Composition Rules}

\subsubsection*{Trivial (per-primitive, derivable by construction)}

\textbf{R-CLAUS-T:} \texttt{clauses(tensor(A,B))[p] = ZFCT\_TEMPLATES[p][max\_ord(A[p], B[p])]}

The clause for each primitive in a tensor product is simply the template for whichever input had the higher ordinal value. Per-primitive independence is total --- no cross-primitive information enters a single clause.

\textbf{R-CLAUS-M:} Same rule using \texttt{min\_ord} for the meet.

\subsubsection*{Non-trivial (cross-primitive, emergent)}

\textbf{R-FROB-BARR:} \\
FROB $\in$ clauses(tensor(A,B)) \textbf{iff} $A[\Phi] = \Phi_{\}}$ \textbf{or} $B[\Phi] = \Phi_{\}}$

$\Phi_{\}}$ (ord 4) is the unique maximum of the $\Phi$ dimension. Therefore FROB cannot be synthesized from two inputs both having $\Phi <$ ord 4. This is the algebraic expression of the \textbf{Frobenius cliff}. Zero violations observed across 600 operations.

\textbf{R-FROB-CONT:} \\
If FROB $\in$ clauses(A), then FROB $\in$ clauses(tensor(A, X)) for all X. \\
(Frobenius is forward-preserved.)

\textbf{R-FIXPT-T:} \\
FIXPT $\in$ clauses(tensor(A,B)) iff $A[\hat{\varphi}] \geq$ ord 1 \textbf{or} $B[\hat{\varphi}] \geq$ ord 1

\textbf{R-TIER-EMRG:} \\
tier(tensor(A,B)) can strictly exceed max(tier(A), tier(B))

The mechanism is exclusively cross-primitive: FROB lives in clause[$\Phi$]; FIXPT lives in clause[$\hat{\varphi}$]. \\
Observed rate: \textbf{6.3\%} of 300 random tensor pairs.

\textbf{R-ZFCT-ABS:} \\
ZFCt is the absorption element for all six promotion channels. \\
\texttt{tensor(X, ZFCt)[c] =} ZFCt value for each promoted primitive c, $\forall$ X.

\begin{center}
\rule{0.5\textwidth}{0.4pt}
\end{center}

\subsection*{Statistical Results}

\begin{table}[htbp]
\centering
\small

\begin{tabularx}{\textwidth}{>{\centering\arraybackslash}X >{\centering\arraybackslash}X}
\toprule
\textbf{Metric} & \textbf{Value} \\
\midrule

Catalog entries & 2706 (2697 valid) \\
Pairs scanned & 300 tensor + 300 meet \\
Total observations & 600 \\
Tensor tier emergences & 19 / 300 (6.3\%) \\
FROB synthesized ex nihilo & 0 / 600 (0\%) \\
\bottomrule
\end{tabularx}
\end{table}

\begin{center}
\rule{0.5\textwidth}{0.4pt}
\end{center}

\subsection*{Interpretation}

\subsubsection*{The Foundational Split}

\begin{itemize}
    \item \textbf{ZFC} is the ourobourically inert carrier of Frobenius symmetry ($\Phi_{\}}$).
    \item \textbf{ZFCt} is the ourobourically active temporal extension, carrying criticality ($\hat{\varphi}_{\text{ÿ}}$) and the six promotion channels, but operating with a reduced $\Phi$.
\end{itemize}

\subsubsection*{Reunion via Tensor}

$\text{ZFC} \otimes \text{ZFCt} = O_\infty$ is the algebraic statement that the foundation and its temporal extension are \textbf{complementary half-objects}. \\
Each supplies the component the other lacks. Their tensor product crosses the Frobenius cliff.

\subsubsection*{$O_\infty$ in the Wild}

Three entries sit at $O_\infty$ natively:

\begin{itemize}
    \item Temporal Mathematics
    \item Einstein field equations
    \item IUG (Mochizuki's Inter-Universal Teichmüller Theory)
\end{itemize}

They possess both FROB and FIXPT by constitution, not by composition.

\begin{center}
\rule{0.5\textwidth}{0.4pt}
\end{center}

\subsection*{The Algebra in Summary}

\begin{table}[htbp]
\centering
\small

\begin{tabularx}{\textwidth}{>{\raggedright\arraybackslash}X >{\centering\arraybackslash}X >{\centering\arraybackslash}X >{\centering\arraybackslash}X >{\raggedright\arraybackslash}X}
\toprule
\textbf{Structure} & \textbf{FROB?} & \textbf{FIXPT?} & \textbf{Tier} & \textbf{Notes} \\
\midrule

ZFC & yes & no & $O_0$ & Non-critical Frobenius \\
ZFCt & no & yes & $O_2^\dagger$ & Critical, non-Frobenius \\
$\text{ZFC} \otimes \text{ZFCt}$ & yes & yes & $O_\infty$ & Tier emergence \\
$\text{ZFC} \sqcap \text{ZFCt}$ & --- & --- & $O_0$ & Destructive meet \\
Temporal Math & yes & yes & $O_\infty$ & Native \\
Einstein EFE & yes & yes & $O_\infty$ & Native \\
IUG / Mochizuki & yes & yes & $O_\infty$ & Native \\
Navier-Stokes & no & yes & $O_2^\dagger$ & Critical, non-Frobenius \\
Schrödinger & no & yes & $O_2^\dagger$ & Critical, non-Frobenius \\
\bottomrule
\end{tabularx}
\end{table}

The Frobenius barrier is absolute: no tensor operation can produce $\Phi_{\}}$ from two non-FROB inputs. Tier emergence to $O_\infty$ requires exactly the complementary pairing.

\begin{center}
\rule{0.5\textwidth}{0.4pt}
\end{center}

\end{document}
